% =====================================================================
%  CAPÍTULO 3 — ESTADO DEL ARTE
% =====================================================================
% Requiere en el preámbulo: booktabs, array, cleveref, siunitx
% Aquí va qué se ha hecho y por qué, a nivel general. Cómo funciona cada
% técnica (CTC, frontend, LoRA...) está en el marco teórico (capítulo 5).
% =====================================================================

\chapter{Estado del arte}
\label{cap:estado-del-arte}

Este capítulo repasa lo que se ha hecho hasta ahora en la decodificación del habla a partir de la actividad cerebral y en la reutilización de modelos de habla. Primero se comparan las técnicas de registro y se recogen las cifras de las neuroprótesis del habla más recientes (\cref{sec:sota-bci,sec:sota-neuroprotesis}). Después se explica cómo están construidos los decodificadores actuales y por qué (\cref{sec:sota-anatomia}), y qué resultados se obtienen en el \textit{benchmark} Brain-to-Text~'25 (\cref{sec:sota-benchmarks}). Por último se resumen los modelos de habla preentrenados y los trabajos que los han llevado a la señal neuronal (\cref{sec:sota-habla}), se comparan todos los sistemas vistos (\cref{sec:sota-resumen}) y se sitúa este trabajo respecto a ellos (\cref{sec:sota-posicionamiento}). En el capítulo aparecen algunos conceptos, como las métricas de error o las partes de un modelo, que se explican más adelante, en el marco teórico (\cref{cap:marco-teorico}).

\section{Modalidades de registro en interfaces cerebro-computador}
\label{sec:sota-bci}

La técnica de registro condiciona qué información se puede extraer y qué tipo de decodificador tiene sentido construir. Las técnicas utilizadas en BCI del habla van de menos a más invasivas. En general, cuanto más invasiva es la técnica, mejor es la señal:

\begin{description}
\item[Electroencefalografía (EEG).] Es una técnica no invasiva, barata y fácil de conseguir. Registra, a través del cráneo, la actividad sumada de grandes grupos de neuronas, por lo que la señal tiene mucho ruido. Funciona bien para elegir entre pocas opciones, pero la decodificación de texto a partir de EEG está en discusión: se ha visto que algunos modelos publicados rinden casi igual utilizando solo ruido que utilizando EEG \cite{jo2024eegtotext}. Además, suele presentar problemas de portabilidad y comodidad, y es la técnica más afectada por el paso del tiempo.
\item[Estereoelectroencefalografía (sEEG).] Utiliza electrodos de profundidad que se implantan con fines diagnósticos en clínicas, a través de pequeños orificios y sin necesidad de abrir el cráneo \cite{verwoert2022}. Tiene buena resolución temporal y permite llegar a estructuras profundas, pero la colocación de los electrodos responde a criterios clínicos y no a las necesidades del decodificador. Por ello, no suele utilizarse para este tipo de problemas.
\item[Electrocorticografía (ECoG).] Utiliza rejillas de electrodos colocadas sobre la superficie de la corteza, lo que obliga a abrir el cráneo \cite{kohler2022}. Estas rejillas cubren una zona amplia, son estables durante mucho tiempo y han permitido construir los primeros sistemas de habla en tiempo real \cite{moses2021,metzger2023}.
\item[Matrices intracorticales de microelectrodos.] Son matrices de electrodos que penetran en la corteza, normalmente de tipo Utah, y registran la actividad de pequeños grupos de neuronas cercanos a cada electrodo. Es la técnica con mayor riesgo quirúrgico, pero con ella se han conseguido prestaciones propias de una conversación \cite{willett2023,card2024}. Es la técnica utilizada para obtener los datos de este trabajo, cuyas características se describen en el \cref{cap:datos}.
\end{description}

La \cref{tab:sota-modalidades} resume las cuatro.

\begin{table}[htbp]
  \centering
  \caption{Comparación de las modalidades de registro usadas en BCI del habla.}
  \label{tab:sota-modalidades}
  \footnotesize
  \setlength{\tabcolsep}{4pt}
  \renewcommand{\arraystretch}{1.2}
  \begin{tabular}{>{\raggedright\arraybackslash}p{0.14\linewidth}
                  >{\raggedright\arraybackslash}p{0.18\linewidth}
                  >{\raggedright\arraybackslash}p{0.17\linewidth}
                  >{\raggedright\arraybackslash}p{0.2\linewidth}
                  >{\raggedright\arraybackslash}p{0.2\linewidth}}
    \toprule
    \textbf{Modalidad} & \textbf{Cirugía} & \textbf{Resolución espacial} & \textbf{Cobertura} & \textbf{Duración de los registros} \\
    \midrule
    EEG           & No & Baja & Toda la cabeza, solo superficie & Sesiones de laboratorio \\
    sEEG          & Mínima, sin abrir el cráneo \cite{verwoert2022} & Alta, en puntos dispersos & Dispersa, incluye estructuras profundas & Lo que dura el ingreso clínico \\
    ECoG          & Apertura del cráneo \cite{kohler2022} & Media & Una zona amplia de la superficie & Larga: unas 81 semanas en \cite{moses2021} \\
    Intracortical & Apertura del cráneo e implante penetrante & Muy alta: grupos pequeños de neuronas & Pequeña: $3{,}2 \times 3{,}2$~mm por matriz \cite{card2024} & Larga: 20 meses en Brain-to-Text~'25 \cite{b2t25} \\
    \bottomrule
  \end{tabular}
\end{table}

\section{Evolución de las neuroprótesis del habla}
\label{sec:sota-neuroprotesis}

En esta sección se recogen los trabajos que marcaron cada salto en el campo de las neuroprótesis. Es importante entenderlos para poder justificar decisiones más tarde.

Uno de los primeros sistemas que produjo lenguaje de forma continua decodificaba la escritura: el participante intentaba escribir en vez de hablar. Llegó a unos 90 caracteres por minuto, con más de un \SI{99}{\percent} de acierto usando un autocorrector \cite{willett2021}.

Con el habla intentada, Moses \textit{et al.} decodificaron frases en tiempo real con ECoG en una persona con anartria, con un vocabulario de 50 palabras y una tasa de error por palabra (WER) del \SI{25.6}{\percent} \cite{moses2021}. En 2023, dos trabajos pasaron a vocabularios mucho más grandes. Metzger \textit{et al.}, con ECoG de alta densidad y también síntesis de voz y un avatar, llegaron a un \SI{25}{\percent} de WER sobre 1.024 palabras \cite{metzger2023}. Willett \textit{et al.}, con matrices intracorticales, llegaron a un \SI{23.8}{\percent} sobre 125.000 palabras \cite{willett2023}.

El sistema de Card \textit{et al.}, del que proceden los datos de este trabajo, es el referente actual. Con 30 minutos de datos de calibración alcanzó un \SI{99.6}{\percent} de acierto con 50 palabras el primer día, y un \SI{90.2}{\percent} con 125.000 palabras el segundo día, tras 1,4 horas más de datos. Mantuvo un \SI{97.5}{\percent} de acierto durante 8,4 meses, y el participante lo usó para conversar más de 248 horas \cite{card2024}. Su decodificador combina una red recurrente que predice fonemas con un modelo de $n$-gramas y un modelo de lenguaje de gran tamaño que reordena las frases candidatas.

Desde entonces el campo se ha abierto mucho. Los datos de Willett \textit{et al.} y de Card \textit{et al.} se publicaron en los retos Brain-to-Text '24 y '25 \cite{b2t24lessons,b2t25}, y ahora muchos grupos sin acceso a pacientes trabajan sobre los mismos registros e intentan mejorar sus cifras. En el reto de 2025, el sistema de referencia de los organizadores tenía un \SI{6.7}{\percent} de WER, y los primeros puestos bajaron hasta cerca del \SI{1.5}{\percent} juntando varios decodificadores y modelos de lenguaje de gran tamaño \cite{b2t25}. Debemos tener en cuenta que estos resultados se obtienen en diferido, sobre datos ya grabados, y no se han probado en tiempo real con el participante, como sí se hizo con el sistema de Card.

\section{Anatomía de un decodificador \textit{brain-to-text}}
\label{sec:sota-anatomia}

Los sistemas anteriores siguen casi todos el mismo esquema, heredado del reconocimiento de voz clásico, y lo van mejorando: un codificador neuronal convierte la señal en probabilidades de fonemas o palabras y un modelo de lenguaje forma las frases \cite{crossspecies2025}. Esta sección explica por qué se ha llegado a ese esquema y hacia dónde está cambiando.

\subsection{El cambio entre días}

Lo que más condiciona el diseño es que la señal cambia de un día a otro, así que un decodificador entrenado con los datos de un día va perdiendo precisión \cite{sussillo2016}. La respuesta del campo ha ido cambiando. Willett \textit{et al.} añadieron delante de un modelo compartido una capa de entrada propia de cada día \cite{willett2021}, una solución (\cref{sec:capa-sesion}) que se ha mantenido en los sistemas posteriores \cite{willett2023,card2024}. Como cada día nuevo necesita datos etiquetados para su capa, el siguiente paso fue reducir esa calibración: Card \textit{et al.} calibran el sistema con 30 minutos de datos \cite{card2024}, y Fan \textit{et al.} lo recalibran con las propias frases que decodifica, sin pedir datos nuevos al usuario \cite{fan2023corp}. Los sistemas más recientes comparten parámetros entre días cercanos \cite{brainwhisperer}.

\subsection{El codificador}

El codificador de referencia sigue siendo una red recurrente, la de los sistemas de Willett \textit{et al.} y Card \textit{et al.} y la del \textit{benchmark} \cite{willett2023,card2024,b2t25}. Se mantiene porque ha costado mucho mejorarla: en el Brain-to-Text~'24, otras arquitecturas más recientes no consiguieron superarla \cite{b2t24lessons}. Los avances recientes llegan cuando se añade otra fuente de información: preentrenar el codificador con datos neuronales de muchas tareas, personas e incluso especies \cite{crossspecies2025}, partir de un modelo de voz preentrenado \cite{brainwhisperer} o usar arquitecturas de reconocimiento de voz como Conformer \cite{khanday2026}.

\subsection{La salida y el modelo de lenguaje}

La unidad de salida más usada es el fonema. La corteza motora controla la articulación y no la ortografía, y con unos 40 fonemas el problema es más sencillo que con miles de palabras \cite{willett2023}.

Pero los fonemas no son texto. Para pasar a palabras hace falta un diccionario de pronunciación, que indica qué fonemas forman cada palabra, y un modelo de lenguaje. En este paso está buena parte del rendimiento final. Los sistemas de referencia usan un modelo de $n$-gramas, que estima la probabilidad de cada palabra a partir de las anteriores \cite{jurafsky2026slp}, y reordenan sus resultados con modelos de lenguaje grandes.

Este esquema es caro. Por ejemplo, un modelo de 5-gramas necesita unos \SI{300}{\giga\byte} de memoria y cada decodificación tarda unos \SI{750}{\milli\second}, frente a menos de \SI{2}{\giga\byte} y unos \SI{50}{\milli\second} de un sistema de extremo a extremo \cite{brainwhisperer}. Además, el codificador y el modelo de lenguaje se entrenan por separado y no se pueden optimizar juntos \cite{crossspecies2025}. Por eso están apareciendo sistemas de extremo a extremo, que pasan de la señal al texto con un solo modelo, prediciendo directamente letras o palabras \cite{khanday2026,brainwhisperer}. Todavía quedan por detrás de los sistemas con modelo de lenguaje (\cref{sec:sota-benchmarks}), pero son más baratos y sencillos de usar y están aumentando rápidamente su precisión.

\section{El \textit{benchmark} Brain-to-Text~'25}
\label{sec:sota-benchmarks}

Los \textit{benchmarks} públicos, con particiones fijas y una métrica común, han acelerado mucho el campo. El Brain-to-Text~'25, el que se usa en este trabajo, se describe en el \cref{cap:datos}. Su sistema de referencia, la red recurrente con fonemas y un modelo de 5-gramas, consigue un \SI{6.7}{\percent} de WER \cite{brainwhisperer}. Las primeras posiciones se consiguen con modelos de lenguaje y, a menudo, combinando varios modelos, como ya pasó en la edición anterior \cite{b2t24lessons}. BrainWhisperer, por ejemplo, quedó cuarto de 466 equipos con una versión ampliada de su sistema con modelo de 5-gramas \cite{brainwhisperer}.

Estas cifras se consiguen con sistemas grandes y lentos. Combinan varios decodificadores con modelos de lenguaje de gran tamaño, y solo el modelo de lenguaje de referencia de la competición necesita cientos de gigabytes de memoria \cite{b2t25}. Hoy sería difícil llevar algo así a un dispositivo que funcione en tiempo real y en casa del usuario. Aun así, marcan hasta dónde se puede llegar con estos datos y por eso es importante tenerlos en cuenta. 


\section{Modelos de habla preentrenados y transferencia a la señal neuronal}
\label{sec:sota-habla}
\label{sec:sota-transferencia}

Por último, el reconocimiento de voz ha pasado a preentrenar modelos con enormes cantidades de audio y a adaptarlos después a cada tarea. Algunos, como wav2vec~2.0 \cite{baevski2020wav2vec}, aprenden de audio sin transcribir y necesitan después un ajuste supervisado. Otros, como Whisper \cite{radford2023whisper}, se entrenan directamente con cientos de miles de horas de audio transcrito. Whisper es un modelo codificador-decodificador, así que, además de leer el audio, trae un decodificador que ha aprendido cómo son las frases. En \textit{brain-to-text} eso interesa, porque el conocimiento del lenguaje es justo lo que le falta a un codificador entrenado con pocos datos. Su inconveniente para la investigación es que no se conocen sus datos ni todo su entrenamiento. 

% OWSM reproduce este tipo de entrenamiento con datos públicos \cite{peng2023owsm,peng2024owsmv31} y es el modelo del que parte este trabajo (\cref{sec:owsm-teoria}).

Llevar estos modelos a la actividad cerebral se ha probado de varias formas. Algunos trabajos toman del modelo de voz su arquitectura o su objetivo de entrenamiento, pero no los pesos aprendidos con audio. BrainTalker proyecta la señal de ECoG al espacio de wav2vec~2.0 para sintetizar voz con muy pocos datos \cite{braintalker}, y Yuan y Makin entrenan wav2vec con ECoG sin etiquetar y mejoran así un decodificador supervisado, incluso pasando de un paciente a otro \cite{ecogssl2024}.

Otros transfieren los pesos aprendidos con audio. Fiedler \textit{et al.} sustituyen el extractor de audio de wav2vec~2.0 por otro pensado para señal neuronal y comparan, con los datos del participante, el ajuste completo de los pesos preentrenados, el entrenamiento desde cero y el modelo preentrenado congelado. El ajuste completo da el mejor resultado, con un \SI{18.54}{\percent} de error por carácter y un \SI{30.97}{\percent} de WER sin modelo de lenguaje, unos 20 puntos de error por carácter mejor que desde cero \cite{fiedler2025wav2vec2brain}. BrainWhisperer lleva la idea a Whisper y a los datos del Brain-to-Text~'25. Adapta la entrada y ajusta el decodificador de Whisper, entre otros cambios, y llega a un \SI{10.2}{\percent} de WER.

\section{Comparativa de los sistemas revisados}
\label{sec:sota-resumen}

La \cref{tab:sota-resumen} reúne los sistemas que han aparecido en este capítulo. No todas las cifras se pueden comparar entre sí. Los sistemas clínicos se evaluaron en tiempo real, cada uno con su participante y su vocabulario, mientras que el resto se evalúa en diferido. Pero podemos usar estos datos para discutir los resultados de este trabajo.

\begin{table}[H]
  \centering
  \caption{Resumen de los sistemas revisados en este capítulo.}
  \label{tab:sota-resumen}
  \footnotesize
  \setlength{\tabcolsep}{3.5pt}
  \renewcommand{\arraystretch}{1.15}
  \begin{tabular}{>{\raggedright\arraybackslash}p{0.22\linewidth}
                  l l l
                  >{\raggedright\arraybackslash}p{0.14\linewidth}
                  l}
    \toprule
    \textbf{Trabajo} & \textbf{Registro} & \textbf{Participante} & \textbf{Vocabulario} & \textbf{WER} & \textbf{Velocidad} \\
    \midrule
    \multicolumn{6}{l}{\textit{Sistemas clínicos, en tiempo real}} \\
    Moses \textit{et al.} \cite{moses2021}     & ECoG          & Anartria    & 50 palabras      & 25,6\,\%         & 15 ppm \\
    Willett \textit{et al.} \cite{willett2021} & Intracortical & Tetraplejia & Escritura        & $<$\,1\,\%$^{*}$ & 90 cpm \\
    Metzger \textit{et al.} \cite{metzger2023} & ECoG          & Anartria    & 1.024 palabras   & 25\,\%           & 78 ppm \\
    Willett \textit{et al.} \cite{willett2023} & Intracortical & ELA (T12)   & 125.000 palabras & 23,8\,\%         & 62 ppm \\
    Card \textit{et al.} \cite{card2024}       & Intracortical & ELA (T15)   & 125.000 palabras & $\sim$2,5\,\%    & 32 ppm \\
    \midrule
    \multicolumn{6}{l}{\textit{Brain-to-Text~'25 con modelo de lenguaje externo, en diferido}} \\
    Referencia de la competición \cite{b2t25} & Intracortical & ELA (T15) & 125.000 palabras & 6,7\,\%$^{\dagger}$ & -- \\
    BrainWhisperer \cite{brainwhisperer}      & Intracortical & ELA (T15) & 125.000 palabras & 5,9\,\% (5,2\,\%)$^{\dagger}$ & -- \\
    Primeros puestos \cite{b2t25}             & Intracortical & ELA (T15) & --               & $\approx$\,1,5\,\%$^{\dagger}$ & -- \\
    \midrule
    \multicolumn{6}{l}{\textit{Brain-to-Text~'25 de extremo a extremo, sin modelo de lenguaje externo, en diferido}} \\
    Red recurrente de referencia \cite{brainwhisperer} & Intracortical & ELA (T15) & Abierto & 16,0\,\%$^{\dagger}$ & -- \\
    BrainWhisperer \cite{brainwhisperer}       & Intracortical & ELA (T15) & Abierto & 10,2\,\% (8,7\,\%)$^{\dagger}$ & -- \\
    Khanday \textit{et al.} \cite{khanday2026} & Intracortical & ELA (T15) & Abierto & 23,8\,\%$^{\ddagger}$ & -- \\
    \bottomrule
  \end{tabular}

  \vspace{0.4em}
  \parbox{\linewidth}{\raggedright\footnotesize ppm: palabras por minuto; cpm: caracteres por minuto. Entre paréntesis, entrenando también con datos de otros participantes (BrainWhisperer) o combinando varios modelos (BIT). $^{*}$Error por carácter, en diferido y con autocorrector. $^{\dagger}$Partición de prueba de la competición. $^{\ddagger}$Error por carácter en la partición de validación.}
\end{table}

\section{Posicionamiento de este trabajo}
\label{sec:sota-posicionamiento}

De esta revisión salen tres huecos, que son los que aborda este trabajo. Primero, la transferencia desde modelos de habla funciona en algunos casos, pero no se sabe bien qué parte la explica. Fiedler \textit{et al.} usan un modelo que solo tiene codificador y otro participante, y BrainWhisperer cambia muchas cosas a la vez para conseguir la mejor cifra. Falta separar, con un mismo presupuesto, qué aporta cada parte de un modelo codificador-decodificador (el \textit{frontend}, el codificador y el decodificador), y hacerlo con un modelo abierto como OWSM.

Segundo, en el \textit{benchmark} las frases de validación y de prueba son de días que también tienen frases de entrenamiento (\cref{sec:b2t25}), así que las cifras publicadas no dicen cómo se comporta un modelo en un día nuevo. Los trabajos que estudian el paso del tiempo lo hacen con sistemas recurrentes y modelo de lenguaje \cite{fan2023corp}. Khanday \textit{et al.} señala que el cambio entre sesiones es lo que más hace variar su error. Falta ver cómo se comporta en días futuros un sistema de extremo a extremo basado en un modelo de voz, y cuántos datos necesita.

Por último, las primeras posiciones necesitan muchos recursos. Este trabajo se limita a una GPU, un solo modelo y ningún modelo de lenguaje externo, así que su referencia son los sistemas de extremo a extremo de la \cref{tab:sota-resumen} y no la clasificación. Aun así, parte con desventajas conocidas frente a ellos: no usa datos de otros participantes y su presupuesto de cómputo es pequeño. Son técnicas que se pueden añadir después sobre el mismo sistema, por tanto la distancia mide lo que queda por ganar.

El trabajo más cercano es el de Khanday \textit{et al.}, que usa los mismos datos y un Conformer entrenado desde cero que predice letras directamente, sin modelo de lenguaje, y obtiene un \SI{23.80}{\percent} de error por carácter en validación \cite{khanday2026}. Este trabajo sigue la misma idea de extremo a extremo, pero parte de un modelo de voz preentrenado aprovechando su decodificador.
